%%% FIELD def-ledger
For the unequal three-arm contrast, put
\[
\mathsf B=[v,d^\dagger,\mathbf 1]
=\begin{pmatrix}
1/2&0&1\\
-1/2&3/4&1\\
-1/2&-1/4&1
\end{pmatrix},
\qquad
\mathsf B^{-1}=
\begin{pmatrix}
1&-1/4&-3/4\\
0&1&-1\\
1/2&1/8&3/8
\end{pmatrix}.
\]
Define the signed response directions
\[
\begin{aligned}
s_x^+&=\lambda_v^+-3p_0^\dagger,&
s_x^-&=3p_0^\dagger-\lambda_v^-,\\
s_h^+&=\lambda_d^+-2p_0^\dagger,&
s_h^-&=2p_0^\dagger-\lambda_d^-,\\
\gamma_0&=\delta_{111}-\delta_{000}.
\end{aligned}
\]
For an allocation proportion \(q\) in the open simplex, set
\[
H(q)=\frac1{q_1}+\frac1{16q_2}+\frac9{16q_3},
\qquad
z_a(q)=\frac{c_a^\dagger}{q_aH(q)},
\]
and define \(\beta(q)\) and \(\gamma(q)\) by
\[
\begin{pmatrix}1\\ \beta(q)\\ \gamma(q)\end{pmatrix}
=\mathsf B^{-1}z(q).
\]
Equivalently,
\[
\beta(q)=\frac{3/q_3-1/q_2}{4H(q)},
\qquad
\gamma(q)=\frac{16/q_1-1/q_2-9/q_3}{32H(q)}.
\]
For a verified spline record define
\[
J_x(k)=\frac{4\int_{\mathbb R}|(\varphi_k^\dagger)'|^2}
{\int_{\mathbb R}(\varphi_k^\dagger)^2},
\quad
J_{\max}(k,\varepsilon)=J_x(k)+\frac{325}{2\varepsilon^2},
\quad
L(k,\varepsilon)=\frac{T_k^\dagger}{2}+\frac{7\varepsilon}{4}.
\]
The following finite maximum explicitly enumerates the eight response types and the six arm--outcome cells:
\[
\begin{aligned}
C(k,\varepsilon)=1+
\max_{(e_x,e_h,e_z)\in\{-1,1\}^3}
\Bigg[&\max_{t\in\{0,1\}^3}
\left|e_xT_k^\dagger s_{x,t}^{e_x}
+e_h\varepsilon s_{h,t}^{e_h}
+e_z\varepsilon\gamma_{0,t}\right|\\
&+\max_{\substack{a\in\{1,2,3\}\\ b\in\{0,1\}}}
2\left|(2b-1)
\{e_xT_k^\dagger v_a+e_h\varepsilon d_a^\dagger
+e_z\varepsilon\}\right|\Bigg],
\end{aligned}
\]
where \(s^1=s^+\) and \(s^{-1}=s^-\).  Put
\[
\Gamma(k,\varepsilon)=4C(k,\varepsilon)+\frac{L(k,\varepsilon)}6,
\qquad
\mathcal W(k,\varepsilon)=\frac58T_k^\dagger+\frac3{16}\varepsilon.
\]
For rational \(\eta>0\), let \(N(k,\varepsilon,\eta)\) be the least positive integer \(N\) satisfying all of the following decidable inequalities:
\[
\begin{gathered}
N^{-1/3}(3T_k^\dagger+4\varepsilon)\le\frac12,
\qquad
4N^{-2/3}L(k,\varepsilon)^2\le\frac12,\\
8N^{-1/3}\Gamma(k,\varepsilon)^2
+12\Gamma(k,\varepsilon)N^{-11/72}
+32N^{-2/3}\Gamma(k,\varepsilon)L(k,\varepsilon)^2
+8N^{-1/3}L(k,\varepsilon)^2\le\eta,\\
\{4+4N^{-11/36}\}J_{\max}(k,\varepsilon)N^{-1/36}\le1,
\qquad
8N^{-2/3}\Gamma(k,\varepsilon)^2\le1,\\
16L(k,\varepsilon)^2N^{-13/36}
+4\mathcal W(k,\varepsilon)N^{-1/36}
+64\Gamma(k,\varepsilon)L(k,\varepsilon)^2N^{-25/36}\\
\hspace{28mm}
+16\Gamma(k,\varepsilon)^2
\{1+8L(k,\varepsilon)^2N^{-2/3}\}N^{-13/36}\le1,\\
2C(k,\varepsilon)N^{-1/36}\le\frac12,
\qquad
2J_{\max}(k,\varepsilon)N^{-1/36}\le\frac12,
\qquad
2L(k,\varepsilon)^2N^{-13/36}\le\frac12,\\
J_{\max}(k,\varepsilon)N^{-1/3}\le1,
\qquad N^{-1/6}\le\eta,
\qquad 2^{-N}\le\eta.
\end{gathered}
\]
The defining search scans positive integers in increasing order, making every comparison by rational arithmetic and isolating intervals for the finitely many positive algebraic roots.

%%% FIELD def-primal-ledger
At \(c=c^\dagger\), run the rational interval Airy oracle until it returns a positive rational lower endpoint and a rational upper endpoint
\(\Lambda_{c^\dagger}<\overline\Lambda^\dagger\).  Choose an integer
\(R\) satisfying
\[
\frac58R-\overline\Lambda^\dagger\ge1,
\qquad
\frac58R\ge2\overline\Lambda^\dagger.
\]
On \([-R,R]\), validated Picard integration and interval subdivision are refined, and the first rational candidates
\(\overline H_0^\dagger\), \(\overline H_1^\dagger\),
\(\overline H_2^\dagger\), and \(\overline D^\dagger\) passing the resulting interval inequalities for
\[
\|h_{c^\dagger}'\|_\infty,
\quad \operatorname{Lip}(h_{c^\dagger}'),
\quad \operatorname{Lip}(h_{c^\dagger}^2),
\quad \sup_{x\ne0}\frac{|xh_{c^\dagger}'(x)|}{\sqrt{|x|}},
\]
and outside that interval uses the rational Riccati verifier with
\(q=(5/8)|x|-\Lambda_{c^\dagger}\), barriers
\(\sqrt q\) and \(\sqrt q+(5/8)/q\), and the differentiated Riccati identities.
are selected.  Set
\[
N_0^\dagger=
\max\left\{1,
\left\lceil
\left(\frac{2\overline H_0^\dagger}{5/8}\right)^3
\right\rceil\right\},
\]
\[
B_{\mathrm{eff}}^\dagger=
5\left(\frac65\right)^2\overline H_1^\dagger
+\frac65\overline H_2^\dagger+2\overline D^\dagger
+4(N_0^\dagger)^2+\overline\Lambda^\dagger N_0^\dagger.
\]
These formulas define the seven displayed outputs whenever the interval search returns.

%%% FIELD stmt-van-trees
Fix integers \(n,k\ge1\), rational \(\varepsilon>0\), and an allocation count
\(r\in\mathcal R_{n,3}\).  Suppose
\[
\alpha_n(3T_k^\dagger+4\varepsilon)\le\frac12,
\qquad
4\alpha_n^2L(k,\varepsilon)^2\le\frac12.
\]
Put \(q=r/n\).  If every \(q_a>0\) and
\(\|q-q^\dagger\|_\infty\le n^{-11/72}\), then
\[
\begin{aligned}
\mathcal B_{n,k,\varepsilon}^{\dagger,\mathrm A}(r)
\ge\frac1n-n^{-4/3}\Bigg[&J_x(k)+\frac58
\int_{\mathbb R}|x|w_k^\dagger(x)\,dx+\frac{15}{256}\varepsilon\\
&+8\alpha_n\Gamma(k,\varepsilon)^2
+12\Gamma(k,\varepsilon)n^{-11/72}
+32\alpha_n^2\Gamma(k,\varepsilon)L(k,\varepsilon)^2
+8\alpha_nL(k,\varepsilon)^2\Bigg].
\end{aligned}
\tag{L}
\]
If every \(q_a>0\), \(\|q-q^\dagger\|_\infty>n^{-11/72}\), and the exterior and zero-face inequalities in the definition of
\(N(k,\varepsilon,\eta)\) hold at \(n\), then
\[
\mathcal B_{n,k,\varepsilon}^{\dagger,\mathrm A}(r)\ge n^{-1}.
\tag{X}
\]
If some \(r_a=0\) and
\(J_{\max}(k,\varepsilon)\le n^{1/3}\), then
\[
\mathcal B_{n,k,\varepsilon}^{\dagger,\mathrm A}(r)\ge n^{-1}.
\tag{Z}
\]
These are lower bounds for the posterior risk of the exact random finite-population target, not for its conditional mean.

%%% FIELD proof-van-trees
Write \(\lambda=(x,h,\zeta)\) for the three prior coordinates and
\(p=p_{n,x,h,\zeta}^{\dagger,\mathrm A}\).  Conditional on \(\lambda\), the multinomial draw followed by a uniform labeling within the response-count orbit is exactly the law of independent labeled response types \(t_1,\ldots,t_n\) with common law \(p\).  This follows by multiplying the multinomial mass of a count vector by the reciprocal of the number of its labelings.  The fixed allocation labels may therefore be conditioned upon without changing conditional independence across units.

Define
\[
M=\frac1n\sum_{i=1}^n\theta_{t_i},\qquad
\mu=\mathbb E_\lambda\theta_t=\alpha_nx,\qquad
\mu_a=\mathbb E_\lambda b_{t,a},
\]
\[
V_a=\operatorname{Var}_\lambda(b_{t,a})=\frac14-\mu_a^2,\qquad
C_a=\operatorname{Cov}_\lambda(\theta_t,b_{t,a}).
\tag{1}
\]
The tangent identities give
\(\mu_a=\alpha_n(xv_a+hd_a^\dagger+\zeta)\), whence
\(|\mu_a|\le\alpha_nL(k,\varepsilon)\).  The second premise implies
\(V_a\ge1/8>0\), so every division below is valid.  Since \(b_{t,a}\) has only the values \(-1/2\) and \(1/2\), its conditional regression is affine and hence exact:
\[
\mathbb E_\lambda(\theta_t\mid b_{t,a})
=\mu+\frac{C_a}{V_a}(b_{t,a}-\mu_a).
\tag{2}
\]
If \(O\) denotes all observed labels and outcomes, conditional independence and \((2)\) give
\[
m_\lambda(O):=\mathbb E(M\mid\lambda,O)
=\mu+\frac1n\sum_{i=1}^n
\frac{C_{A_i}}{V_{A_i}}(b_{t_i,A_i}-\mu_{A_i}).
\tag{3}
\]
The conditional-variance identity for the nested sigma-fields generated by \(O\) and by \((\lambda,O)\) is
\[
\mathbb E\operatorname{Var}(M\mid O)
=\mathbb E\operatorname{Var}(M\mid\lambda,O)
+\mathbb E\operatorname{Var}(m_\lambda(O)\mid O).
\tag{4}
\]
Using \((2)\) separately for every observed arm gives the exact first term
\[
\mathbb E\operatorname{Var}(M\mid\lambda,O)
=\mathbb E_\lambda\left[
\frac1{n^2}\sum_{a=1}^3r_a
\left\{\operatorname{Var}_\lambda(\theta_t)-\frac{C_a^2}{V_a}\right\}
\right].
\tag{5}
\]

Let \(D\) be any one of the ledger directions for which
\(D\mu=\alpha_n\) and \(D\mu_a=\alpha_nz_a\).  Directional integration by parts is legitimate for the compact product density by the established tangent-score identity; its weak score is denoted by \(S_D\), and its prior information by \(J_D=\mathbb E S_D^2\).  The score of the joint law of \((\lambda,O)\) is
\[
\mathcal S_D=S_D(\lambda)+
\sum_{i=1}^n
\frac{\alpha_nz_{A_i}(b_{t_i,A_i}-\mu_{A_i})}{V_{A_i}}.
\tag{6}
\]
The posterior mean of \(\mathcal S_D\) given \(O\) is zero.  Applying integration by parts to \(m_\lambda(O)\), and then using \((3)\), yields
\[
\begin{aligned}
\mathbb E[(\mathbb E[m_\lambda(O)\mid O]-m_\lambda(O))\mathcal S_D]
&=\mathbb E[D m_\lambda(O)]\\
&=\alpha_n\left[1-sum_aq_a
\mathbb E_\lambda\left(\frac{z_aC_a}{V_a}\right)\right]
=:A_D.
\end{aligned}
\tag{7}
\]
The cross term between prior and likelihood scores has conditional expectation zero, so
\[
\mathbb E\mathcal S_D^2
=J_D+\alpha_n^2\sum_ar_a
\mathbb E_\lambda\left(\frac{z_a^2}{V_a}\right)=:J_D+I_D.
\tag{8}
\]
Cauchy--Schwarz is applied only now, to \((7)\)--\((8)\), and gives
\[
\mathbb E\operatorname{Var}(m_\lambda(O)\mid O)
\ge\frac{A_D^2}{J_D+I_D}.
\tag{9}
\]
Thus \((4)\), \((5)\), and \((9)\) are the requested generalized van Trees inequality for the exact target.

We record finite bounds used in all three strata.  On each sign cell put
\(\ell_a=xv_a+hd_a^\dagger+\zeta\).  The finite maximum defining
\(C(k,\varepsilon)\) bounds every response-probability perturbation coefficient.  Since \(|\theta_tb_{t,a}|\le1/2\) and there are eight types,
\[
\left|\sum_t\theta_tb_{t,a}
\{xs_{x,t}^{\operatorname{sign}(x)}
+hs_{h,t}^{\operatorname{sign}(h)}+\zeta\gamma_{0,t}\}\right|
\le4C(k,\varepsilon).
\]
The first positivity premise gives \(\alpha_nT_k^\dagger\le1/6\).  Since the raw moment at \(p_0^\dagger\) is zero, it follows that
\[
|C_a|\le\alpha_n\Gamma(k,\varepsilon),
\qquad
\sum_aq_a\mathbb E\frac{C_a^2}{V_a}
\le8\alpha_n^2\Gamma(k,\varepsilon)^2.
\tag{10}
\]
The exact eight-type calculation is
\[
\sigma_\lambda^2:=\operatorname{Var}_\lambda(\theta_t)
=\alpha_n\left(\frac58|x|+\frac3{16}|h|\right)-\alpha_n^2x^2,
\tag{11}
\]
and hence
\(0\le\sigma_\lambda^2\le\alpha_n\mathcal W(k,\varepsilon)\).

For the local set take \(D=\partial_x\), so \(z=v\), \(J_D=J_x(k)\), and \(v_a^2=1/4\).  Put \(d_q=\|q-q^\dagger\|_\infty\).  The identities
\(4q_a^\dagger v_a=c_a^\dagger\) and
\(\sum_ac_a^\dagger C_a=\sigma_\lambda^2\), together with
\[
0\le \frac1{V_a}-4
=\frac{16\mu_a^2}{1-4\mu_a^2}
\le32\alpha_n^2L(k,\varepsilon)^2,
\tag{12}
\]
give, after averaging over \(\lambda\),
\[
\left|\sum_aq_a\mathbb E\frac{v_aC_a}{V_a}
-\mathbb E\sigma_\lambda^2\right|
\le6d_q\alpha_n\Gamma(k,\varepsilon)
+16\alpha_n^3\Gamma(k,\varepsilon)L(k,\varepsilon)^2.
\tag{13}
\]
Also
\[
\sum_aq_a\mathbb E\frac{v_a^2}{V_a}=1+\rho,
\qquad 0\le\rho\le8\alpha_n^2L(k,\varepsilon)^2.
\tag{14}
\]
For all real \(b\) and \(t\ge0\), the exact square identity
\[
\frac{(1-b)^2}{1+t}-(1-2b-t)=\frac{(b+t)^2}{1+t}\ge0
\tag{15}
\]
applied to \((9)\), with
\(t=\alpha_nJ_x(k)+\rho\), combines with \((5)\) to give
\[
\mathcal B_{n,k,\varepsilon}^{\dagger,\mathrm A}(r)
\ge\frac1n-\frac1n\left[
\mathbb E\sigma_\lambda^2+8\alpha_n^2\Gamma^2
+12d_q\alpha_n\Gamma+32\alpha_n^3\Gamma L^2
+\alpha_nJ_x(k)+8\alpha_n^2L^2\right],
\tag{16}
\]
where arguments \((k,\varepsilon)\) are suppressed only in this display.  Substitute \((11)\), use
\(\mathbb E|H^\dagger|=5\varepsilon/16\), and use
\(d_q\le n^{-11/72}\).  Since \(n^{-4/3}=\alpha_n/n\), equation \((16)\) is exactly \((\mathrm L)\).

For a positive-count exterior allocation take
\(D_q=\partial_x+\beta(q)\partial_h+\gamma(q)\partial_\zeta\).
The established ledger gives \(J_D\le J_{\max}(k,\varepsilon)\), and its exact identities give
\[
\sum_aq_az_aC_a=\frac{\sigma_\lambda^2}{H(q)},
\qquad
\sum_aq_az_a^2=\frac1{H(q)}.
\tag{17}
\]
Set
\[
j=\alpha_nJ_{\max}(k,\varepsilon),
\quad u=8\alpha_n^2L(k,\varepsilon)^2,
\quad \Delta=H(q)-4.
\]
Equations \((10)\), \((12)\), and \((17)\) imply
\[
\left|\sum_aq_a\mathbb E\frac{z_aC_a}{V_a}
-\frac{4\mathbb E\sigma_\lambda^2}{H(q)}\right|
\le\frac{64\alpha_n^3\Gamma(k,\varepsilon)L(k,\varepsilon)^2}{H(q)},
\tag{18}
\]
\[
I_D\le\frac{4n\alpha_n^2(1+u)}{H(q)}.
\tag{19}
\]
Moreover Cauchy--Schwarz on the three allocation coordinates gives the exact identity and bound
\[
H(q)-4=4\sum_{a=1}^3\frac{(q_a-q_a^\dagger)^2}{q_a}
\ge4\|q-q^\dagger\|_\infty^2,
\tag{20}
\]
so \(\Delta\ge\Delta_0:=4n^{-11/36}\) on the exterior.

Let \(b\) denote the sum in \((18)\) and let
\(D_0=4(1+u)+jH(q)\).  Equations \((5)\), \((9)\), and \((19)\) imply
\[
n\mathcal B_{n,k,\varepsilon}^{\dagger,\mathrm A}(r)
\ge \mathbb E\sigma_\lambda^2
-8\alpha_n^2\Gamma^2+\frac{H(q)(1-b)^2}{D_0}.
\tag{21}
\]
After subtracting one and multiplying by \(D_0>0\), discard the nonnegative term
\(D_0\mathbb E\sigma_\lambda^2\) and use \((18)\).  This gives the explicit lower bound
\[
\begin{aligned}
D_0\{n\mathcal B(r)-1\}
\ge{}&\Delta-jH(q)-4u-8\alpha_n\mathcal W
-128\alpha_n^3\Gamma L^2\\
&-32\alpha_n^2\Gamma^2(1+u)
-8\alpha_n^2\Gamma^2jH(q).
\end{aligned}
\tag{22}
\]
The first exterior inequality in the revised threshold gives
\(jH(q)\le\Delta/4\), because \(\Delta/H(q)\) is increasing in \(\Delta\) and \(\Delta\ge\Delta_0\).  The inequality
\(8\alpha_n^2\Gamma^2\le1\) then makes the last term in \((22)\) at most \(\Delta/4\).  The remaining exterior inequality is precisely
\[
4u+8\alpha_n\mathcal W+128\alpha_n^3\Gamma L^2
+32\alpha_n^2\Gamma^2(1+u)\le\Delta_0/2.
\]
Consequently the right side of \((22)\) is nonnegative, proving \((\mathrm X)\).  Dividing this last inequality by
\(2n^{-11/36}\) displays the required decaying factors
\(\mathcal Wn^{-1/36}\), \(\Gamma^2n^{-13/36}\), and
\(L^2n^{-13/36}\); the additional explicit tests involving
\(C n^{-1/36}\), \(J_{\max}n^{-1/36}\), and
\(L^2n^{-13/36}\) reserve separate half-budgets and are monotone.

Finally suppose arm \(a\) is missing.  Use the corresponding live face direction with parameter coordinates
\((1,0,1/2)\), \((1,-4,-1/2)\), or \((1,4/3,-1/2)\).
The tangent-score lemma says that every revealed arm marginal has directional derivative zero.  Thus the likelihood part of \((6)\) vanishes, while \(D\mu=\alpha_n\); equation \((7)\) has \(A_D=\alpha_n\), and \(J_D\le J_{\max}(k,\varepsilon)\).  The first term in \((4)\) is nonnegative, so \((9)\) gives
\[
\mathcal B_{n,k,\varepsilon}^{\dagger,\mathrm A}(r)
\ge\frac{\alpha_n^2}{J_D}
\ge\frac{\alpha_n^2}{J_{\max}(k,\varepsilon)}\ge\frac1n,
\]
where the last step is exactly the stated zero-face test.  This proves \((\mathrm Z)\) and completes the proof.

%%% FIELD stmt-majorants
The certified computation in
\(\mathrm{def:effective\text{-}cdagger\text{-}primal\text{-}ledger}\) terminates and returns rational numbers satisfying
\[
\|h_{c^\dagger}'\|_\infty\le\overline H_0^\dagger,
\quad
\operatorname{Lip}(h_{c^\dagger}')\le\overline H_1^\dagger,
\quad
\operatorname{Lip}(h_{c^\dagger}^2)\le\overline H_2^\dagger,
\quad
|zh_{c^\dagger}'(z)|\le\overline D^\dagger\sqrt{|z|}
\]
for every real \(z\).  In particular, \(B_{\mathrm{eff}}^\dagger\) is a computable positive rational number.

%%% FIELD proof-majorants
Put \(\kappa=5/8\), \(\Lambda=\Lambda_{c^\dagger}\), and
\(h=h_{c^\dagger}\).  The Airy equation and the Riccati identity give, on \(x>0\),
\[
h'(x)=\frac{h(x)^2-(\kappa x-\Lambda)}2.
\tag{1}
\]
Let \(q(x)=\kappa x-\Lambda\).  On an interval where \(q>0\), put
\(p=\sqrt q\), \(u=\sqrt q+\kappa/q\), and let \(f\) and \(g\) be the positive functions with logarithmic derivatives
\(-2f'/f=p\) and \(-2g'/g=u\) that vanish at infinity.  Direct differentiation gives
\[
f''\le qf/4,\qquad g''\ge qg/4.
\]
If \(\Phi\) is the positive square-integrable Airy ground state, the Wronskians
\(\Phi'f-\Phi f'\) and \(\Phi'g-\Phi g'\) have derivatives of opposite signs and vanish at infinity.  Integrating those derivative inequalities from \(x\) to infinity yields
\[
\sqrt q\le h(x)\le\sqrt q+\frac{\kappa}{q}.
\tag{2}
\]
The vanishing of the Wronskians follows from the square-integrable decaying branch and the explicit decays of \(f\) and \(g\); alternatively it follows first along a sequence tending to infinity and then from monotonicity.  Thus \((2)\) is an elementary consequence of the Airy ODE and not an assumed asymptotic expansion.

Equations \((1)\)--\((2)\) imply, whenever \(q\ge1\),
\[
0\le h'(x)\le\frac{\kappa}{\sqrt q}+\frac{\kappa^2}{2q^2}.
\tag{3}
\]
Differentiating \((1)\) gives
\(h''=hh'-\kappa/2\), while
\((h^2)'=2hh'\).  Combining these identities with \((2)\)--\((3)\) yields the explicit exterior bounds
\[
|h''(x)|\le\frac{3\kappa}{2}+\frac{3\kappa^2}{2}+\frac{\kappa^3}{2},
\qquad
|(h^2)'(x)|\le2\kappa+3\kappa^2+\kappa^3.
\tag{4}
\]
When additionally \(\kappa x\ge2\Lambda\), one has \(q\ge\kappa x/2\), and \((3)\) gives
\[
x|h'(x)|\le\{\sqrt{2\kappa}+\kappa\}\sqrt x.
\tag{5}
\]
Evenness of the ground state makes \(h'\) and \(h^2\) even and transfers \((3)\)--\((5)\) to the negative half-line.

The rational Airy oracle supplies a rational upper bound
\(\overline\Lambda^\dagger>\Lambda\), so the two displayed tests defining \(R\) ensure that all exterior premises above hold.  On the compact interval \([-R,R]\), the ground state is strictly positive.  Validated Picard integration therefore produces, after sufficiently fine rational subdivision, a strictly positive rational lower enclosure for the ground state and rational enclosures for it and its derivative.  Substitution into \((1)\),
\(h''=hh'-\kappa/2\), and \((h^2)'=2hh'\) gives certified compact bounds.  Interval widths converge to zero under mesh and Picard-order refinement, so this finite search terminates.  Taking rational maxima of the compact bounds and \((3)\)--\((5)\) produces the four claimed majorants.  All subsequent arithmetic in the definition of \(N_0^\dagger\) and \(B_{\mathrm{eff}}^\dagger\) is finite rational arithmetic.  Hence the entire computation terminates and certifies the asserted global inequalities.

%%% FIELD stmt-eff-qstar
For \(c=c^\dagger\), the total estimator \(\delta_n^\star\) under the independent allocation \(q^\dagger\) satisfies, for every \(n\ge1\),
\[
\max_{y\in\{0,1\}^{n\times3}}
\mathbb E\left[(\delta_n^\star-\tau_{c^\dagger}(y))^2\right]
\le\frac1n-\Lambda_{c^\dagger}n^{-4/3}
+B_{\mathrm{eff}}^\dagger n^{-3/2}.
\]
The displayed remainder constant is the computable rational output of
\(\mathrm{def:effective\text{-}cdagger\text{-}primal\text{-}ledger}\), not an existential analytic constant.

%%% FIELD proof-eff-qstar
Projection onto \([-1,1]\) cannot increase squared loss because
\(\tau_{c^\dagger}(y)\in[-1,1]\).  Analyze the unprojected rule.  For a fixed schedule define
\[
\theta_i=(c^\dagger)^\top(Y_i-\mathbf1/2),
\quad \gamma_i=(g^\star)^\top(Y_i-\mathbf1/2),
\quad X=\sum_i(T_i-\theta_i),
\quad W=\sum_i(U_i-\gamma_i).
\]
The exact-primitives proposition gives \(g^\star=(1,-1/10,-9/10)\), and direct inspection of the three active arms gives
\[
|T_i|=1,\qquad |U_i|\le\frac65,\qquad |\gamma_i|\le1.
\tag{1}
\]
As in the independent score calculation,
\[
\mathbb ET_i=\theta_i,quad \mathbb EU_i=\gamma_i,quad
\mathbb ET_i^2=1,quad \mathbb E(T_iU_i)=1.
\tag{2}
\]
Let \(G=6/5\).  Independence, centering, and \((1)\)--\((2)\) imply
\[
\mathbb EX^2\le n,\qquad \mathbb EW^2\le G^2n,\qquad
\mathbb EW^4\le19G^4n^2,
\tag{3}
\]
because in the fourth-power expansion only one-index fourth moments and two-index products survive, with
\(|U_i-\gamma_i|\le2G\).  Therefore
\[
\mathbb E(|X|W^2)\le5G^2n^{3/2},\qquad
\mathbb E|W|\le G\sqrt n.
\tag{4}
\]

Put \(a_n=n^{-2/3}\),
\(u=a_n\sum_i\gamma_i\), \(Q=\sum_i\theta_i^2\),
\(P=\sum_i\theta_i\gamma_i\), and \(r=Q/n\).
Taylor's formula and the certified Lipschitz majorants give
\[
h(u+a_nW)=h(u)+a_nh'(u)W+R,
\qquad |R|\le\frac12\overline H_1^\dagger a_n^2W^2,
\tag{5}
\]
and
\[
|h(u+a_nW)^2-h(u)^2|
\le\overline H_2^\dagger a_n|W|.
\tag{6}
\]
Using \(\mathbb E(XW)=n-P\), substitution of \((4)\)--\((6)\) into the exact squared-error expansion yields
\[
\begin{aligned}
\mathbb E(\delta_{n,\mathrm{unproj}}^\star-\tau)^2
\le{}&\frac1n-\frac rn
+n^{-4/3}\{h(u)^2-2h'(u)(1-P/n)\}\\
&+\{5G^2\overline H_1^\dagger
+G\overline H_2^\dagger\}n^{-3/2}.
\end{aligned}
\tag{7}
\]
No asymptotic remainder notation occurs in \((7)\).

The dual inequalities imply
\[
r\ge\frac58n^{-1/3}|u|,
\qquad |P|/n\le\frac85r.
\tag{8}
\]
Writing \(d=r-(5/8)n^{-1/3}|u|\ge0\) and using
\(h^2-2h'=(5/8)|u|-\Lambda_{c^\dagger}\), the excess in \((7)\) above
\(n^{-1}-\Lambda_{c^\dagger}n^{-4/3}\) is at most
\[
-\frac dn+\frac{2\overline H_0^\dagger}{5/8}d,n^{-4/3}
+2|uh'(u)|n^{-5/3}
+\{5G^2\overline H_1^\dagger+G\overline H_2^\dagger\}n^{-3/2}.
\tag{9}
\]
For \(n\ge N_0^\dagger\), the first two terms are nonpositive by the definition of \(N_0^\dagger\).  Equation \((1)\) gives \(|u|\le n^{1/3}\), while the certified exterior/compact majorant gives
\[
2|uh'(u)|n^{-5/3}\le2\overline D^\dagger n^{-3/2}.
\tag{10}
\]
Thus the first three summands defining \(B_{\mathrm{eff}}^\dagger\) dominate the remainder for \(n\ge N_0^\dagger\).

For \(1\le n<N_0^\dagger\), the projected estimator and target both lie in \([-1,1]\), so risk is at most four.  Since
\(n^{3/2}\le(N_0^\dagger)^2\) and
\(\Lambda_{c^\dagger}n^{1/6}\le\overline\Lambda^\dagger N_0^\dagger\), the last two summands defining
\(B_{\mathrm{eff}}^\dagger\) dominate the finite prefix after multiplying the desired inequality by \(n^{3/2}\).  The asserted bound follows for every \(n\).

%%% FIELD stmt-eff-grid
Let \(U_n\) be the optimum returned by the rational action-grid program in
\(\mathrm{def:evaluator\text{-}handle}\), with
\(G=\lceil n^{3/4}\rceil\), under the fixed independent allocation
\(q^\dagger\).  Then
\[
R_n(c^\dagger)\le U_n
\le\frac1n-\Lambda_{c^\dagger}n^{-4/3}
+B_{\mathrm{eff}}^\dagger n^{-3/2}+\frac1{4G^2}.
\]
The diagonal evaluator may choose the rational primal error so that
\[
e_{\mathrm{pr},n}\le
\frac{B_{\mathrm{eff}}^\dagger+1/4}
{\lfloor n^{1/7}\rfloor}+2^{-n}.
\]
The right-hand side is a computable rational-valued nonincreasing function of \(n\) tending to zero.

%%% FIELD proof-eff-grid
For each observation \(\omega\), let \(d_\omega\in[-1,1]\) be the action of
\(\delta_n^\star\).  Randomize between the two adjacent grid points whose interval contains \(d_\omega\), with probabilities chosen to preserve conditional mean.  If their spacing is \(G^{-1}\), the randomized action
\(\widetilde d_\omega\) satisfies
\[
\mathbb E(\widetilde d_\omega\mid\omega)=d_\omega,
\qquad
\operatorname{Var}(\widetilde d_\omega\mid\omega)\le\frac1{4G^2}.
\tag{1}
\]
Conditional square completion therefore gives, for every
\(\tau\in[-1,1]\),
\[
\mathbb E\{(\widetilde d_\omega-\tau)^2\mid\omega\}
\le(d_\omega-\tau)^2+\frac1{4G^2}.
\tag{2}
\]
These probabilities are a feasible real point of the rational finite-dimensional linear program.  A rational optimal basic feasible point exists, so program optimality, \((2)\), and the effective fixed-design remainder lemma give
\[
U_n\le n^{-1}-\Lambda_{c^\dagger}n^{-4/3}
+B_{\mathrm{eff}}^\dagger n^{-3/2}+\frac1{4G^2}.
\tag{3}
\]
The returned kernel and independent assignment are feasible in the unrestricted risk game, proving \(R_n(c^\dagger)\le U_n\).

Because \(G\ge n^{3/4}\), multiplication of \((3)\) by \(n^{4/3}\) shows that the positive primal deficit is at most
\((B_{\mathrm{eff}}^\dagger+1/4)n^{-1/6}\), apart from the evaluator's Airy interval width.  The latter is refined below \(2^{-n}\).  For every integer \(n\ge1\),
\[
n^{-1/6}\le n^{-1/7}le\frac1{\lfloor n^{1/7}\rfloor}.
\]
This proves the rational bound in the statement.  Both denominator and dyadic terms are nonincreasing, and each tends to zero, proving the effective tail assertion.

%%% FIELD proof-qstar
Only active arms have positive probability and, on them,
\(q_a^\star=|c_a|/2>0\); hence every displayed ratio in the procedure is defined on its support.  The null-set convention makes the estimator total.  Since the two normalization assumptions imply \(\tau_c(y)\in[-1,1]\), projection onto that interval cannot increase squared loss.

Fix a schedule and put
\[
\theta_i=c^\top(Y_i-\mathbf1/2),\qquad
\gamma_i=(g^\star)^\top(Y_i-\mathbf1/2),
\qquad G_c=\max_{a\in S_c}\left|\frac{g_a^\star}{c_a}\right|.
\]
The support is finite and every denominator is nonzero, so \(G_c<\infty\).  Direct expectation over the independent arm label gives
\[
\mathbb ET_i=\theta_i,quad \mathbb EU_i=\gamma_i,quad
\mathbb ET_i^2=1,quad \mathbb E(T_iU_i)=1,
\tag{1}
\]
where the last two identities use \(q_a^\star=|c_a|/2\),
\(\sum_a|c_a|=2\), and \((g^\star)^\top v=1\).  Also
\(|T_i|=1\), \(|U_i|\le G_c\), and \(|\gamma_i|\le G_c\).

Let \(H_0,H_1,H_2,D<\infty\) bound, respectively,
\(\|h_c'\|_\infty\), the Lipschitz constants of \(h_c'\) and \(h_c^2\), and
\(|xh_c'(x)|/\sqrt{|x|}\).  Their finiteness follows from the unit Airy feedback bounds and the fixed rescaling by \(\kappa(c)>0\).  With
\[
X=\sum_i(T_i-\theta_i),\quad W=\sum_i(U_i-\gamma_i),
\quad Q=\sum_i\theta_i^2,\quad P=\sum_i\theta_i\gamma_i,
\quad r=Q/n,\quad u=n^{-2/3}\sum_i\gamma_i,
\]
independence and expansion of the fourth moment give the explicit bounds
\[
\mathbb EX^2\le n,quad \mathbb EW^2\le G_c^2n,quad
\mathbb EW^4\le19G_c^4n^2,quad
\mathbb E(|X|W^2)\le5G_c^2n^{3/2}.
\tag{2}
\]
Taylor's formula with Lipschitz derivative gives a remainder bounded by
\(H_1n^{-4/3}W^2/2\), while the Lipschitz bound for \(h_c^2\) gives
\(H_2n^{-2/3}|W|\).  Substitution into the exact squared-error identity and use of \((1)\)--\((2)\) therefore give, without asymptotic notation,
\[
\begin{aligned}
\mathbb E(\delta_{n,\mathrm{unproj}}^\star-\tau_c(y))^2
\le{}&\frac1n-\frac rn+n^{-4/3}
\{h_c(u)^2-2h_c'(u)(1-P/n)\}\\
&+(5G_c^2H_1+G_cH_2)n^{-3/2}.
\end{aligned}
\tag{3}
\]
The dual constraint yields
\[
r\ge\kappa(c)n^{-1/3}|u|,\qquad |P|/n\le r/\kappa(c).
\tag{4}
\]
Use the Riccati identity
\(h_c^2-2h_c'=\kappa(c)|u|-\Lambda_c\), and write
\(d=r-\kappa(c)n^{-1/3}|u|\ge0\).  The excess in \((3)\) above the claimed Airy main term is at most
\[
-d/n+2H_0d\,n^{-4/3}/\kappa(c)
+2D\sqrt{|u|}\,n^{-5/3}
+(5G_c^2H_1+G_cH_2)n^{-3/2}.
\tag{5}
\]
Since \(|u|\le G_cn^{1/3}\), the third term is at most
\(2D\sqrt{G_c},n^{-3/2}\).  For
\(n\ge N_c:=\max\{1,\lceil(2H_0/\kappa(c))^3\rceil\}\), the first two terms in \((5)\) have nonpositive sum.  Hence the claimed inequality holds on this tail with the finite displayed constant
\[
5G_c^2H_1+G_cH_2+2D\sqrt{G_c}.
\]
For the finite prefix, projected loss is at most four; enlarging the same constant by
\(4N_c^2+\Lambda_cN_c\) gives a finite \(B_c\) valid for every \(n\ge1\).

The independent procedure is feasible both in the restricted and unrestricted games, so
\[
R_n(c)\le R_n^{\mathrm{ind}}(c)
\le\max_y\mathbb E[(\delta_n^\star-\tau_c(y))^2].
\]
Subtracting from \(n^{-1}\), multiplying by \(n^{4/3}>0\), and taking lower limits proves both concluding inequalities.

%%% FIELD proof-root
Use the corrected fixed-parameter threshold in
\(\mathrm{def:face\text{-}robust\text{-}directional\text{-}ledger}\).  For fixed
\((k,\varepsilon,\eta)\), every coefficient in its finite list is a finite rational number.  Every displayed left-hand side is nonincreasing in the candidate integer, and every term required to vanish is a positive constant times a negative power.  Hence all inequalities eventually hold simultaneously.  Ordered-algebraic isolation decides each comparison, so the increasing search defining \(N(k,\varepsilon,\eta)\) terminates, and every inequality remains valid at later sample sizes.  Put
\[
\overline N_s=\max_{1\le j\le s}N(j,2^{-j},2^{-j}),
\qquad
d(n)=\max\{s\le n:\overline N_s\le n\}
\tag{1}
\]
when the set is nonempty.  Each number in \((1)\) is computable by finite rational and ordered-algebraic comparisons.  For fixed \(s\), every
\(n\ge\max\{s,\overline N_s\}\) makes \(s\) feasible, and feasibility persists.  Hence
\[
k_n=d(n)\uparrow\infty,\qquad
\varepsilon_n^{\mathrm A}=2^{-d(n)}\downarrow0.
\tag{2}
\]
The first two threshold inequalities give
\[
\alpha_n(3T_{k_n}^\dagger+4\varepsilon_n^{\mathrm A})\le\frac12,
\qquad
4\alpha_n^2L(k_n,\varepsilon_n^{\mathrm A})^2\le\frac12.
\tag{3}
\]
Thus the prior fibers are probability laws, every denominator in the exact-target lemma is positive, and the finite-input evaluator terminates.

For \(r\in\mathcal L_n\), apply \((\mathrm L)\) of the exact-target directional van Trees lemma.  The normalized Airy-record inequality and
\(\mathbb E|H^\dagger|=5\varepsilon_n^{\mathrm A}/16\), together with the local-remainder threshold at tolerance \(2^{-k_n}\), give
\[
\mathcal B_{n,k_n,\varepsilon_n^{\mathrm A}}^{\dagger,\mathrm A}(r)
\ge\frac1n-
\left\{\Lambda_{c^\dagger}+2^{1-k_n}
+\frac{15}{256}\varepsilon_n^{\mathrm A}\right\}n^{-4/3}.
\tag{4}
\]
For \(r\in\mathcal X_n\), every exterior threshold inequality is valid by monotonicity, so \((\mathrm X)\) gives
\(\mathcal B(r)\ge n^{-1}\).  For \(r\in\mathcal Z_n\), the zero-face threshold and \((\mathrm Z)\) give the same conclusion.  Consequently conclusion \({\rm(A)}\) holds with the explicit rational errors
\[
e_{\mathrm{loc},n}=2^{1-k_n}+\frac{15}{256}\varepsilon_n^{\mathrm A}+2^{-n},
\qquad e_{\mathrm{shell},n}=e_{\mathrm{zero},n}=0.
\tag{5}
\]
The final summand uses diagonal mode's Airy bracket of width \(2^{-n}\) and satisfies the evaluator's lower-endpoint enclosure convention.

The orbit-Bayes lemma applies by \((3)\).  For every outcome-independent, possibly dependent labeled assignment law and every randomized estimator, it gives
\[
\operatorname{BR}\ge\sum_r\pi(r)\mathcal B(r)
\ge B_{n,k_n,\varepsilon_n^{\mathrm A}}^{\dagger,\mathrm A}.
\tag{6}
\]
The target in \((6)\) is exactly
\(n^{-1}\sum_i\theta_{t_i}=\tau_{c^\dagger}(y)\), by the exact-primitives proposition.  Hence maximum risk dominates \((6)\), and taking the infimum in the unrestricted game proves
\(R_n(c^\dagger)\ge B_{n,k_n,\varepsilon_n^{\mathrm A}}^{\dagger,\mathrm A}\).

On the attaining side, solve the rational action-grid program.  The effective grid lemma gives a total procedure under \(q^\dagger\) with conclusion \({\rm(B)}\), where one may take
\[
e_{\mathrm{pr},n}=
\frac{B_{\mathrm{eff}}^\dagger+1/4}{\lfloor n^{1/7}\rfloor}+2^{-n}.
\tag{7}
\]
This is rational, computable, and tends to zero.

It remains to certify one all-future tail.  Let
\(N_{\mathrm{start}}=\max\{1,N(1,2^{-1},2^{-1})\}\).  For
\(N\ge N_{\mathrm{start}}\), define
\[
E(N)=\max\left\{
\left(2+\frac{15}{256}\right)2^{-d(N)}+2^{-N},
\frac{B_{\mathrm{eff}}^\dagger+1/4}{\lfloor N^{1/7}\rfloor}+2^{-N}
\right\}.
\tag{8}
\]
Extend \(E\) constantly to the finite prefix.  The function is rational-valued, computable, and nonincreasing.  By \((2)\), its first entry tends to zero; the second entry plainly does also.  Equations \((5)\)--\((8)\) prove conclusion \({\rm(C)}\) for every \(n\ge N\).

Finally, by the exact Bayes-risk definition and finiteness of the allocation simplex,
\[
n^{4/3}\{n^{-1}-B_{n,k_n,\varepsilon_n^{\mathrm A}}^{\dagger,\mathrm A}\}
=\max_rG_{n,k_n,\varepsilon_n^{\mathrm A}}^{\dagger,\mathrm A}(r).
\tag{9}
\]
The three strata partition all allocations, so \((4)\)--\((5)\), \((6)\), and the primal bound yield the displayed two-sided risk squeeze in the theorem.  Letting \(N\to\infty\) in \((8)\) proves the stated expansion and
\(\Lambda_{c^\dagger}=C_A(5/8)^{2/3}\in(0,\infty)\).

%%% FIELD proof-sharp
The exact-primitives proposition gives
\(\kappa(c^\dagger)=5/8\) and identifies
\(\tau_{c^\dagger}(m)=\tau_{c^\dagger}(y)\) for every labeled schedule in the response-count orbit.  The Airy definition therefore gives
\[
\Lambda_{c^\dagger}=C_A(5/8)^{2/3}\in(0,\infty).
\tag{1}
\]
For every sufficiently large \(n\), the root theorem supplies the live positivity certificate
\[
\alpha_n(3T_{k_n}^\dagger+4\varepsilon_n^{\mathrm A})\le1,
\tag{2}
\]
the exact-target all-count bounds, and the effective primal procedure.  The exact-target directional van Trees lemma and the orbit-Bayes lemma therefore apply to the random finite-population target itself.  For an arbitrary outcome-independent design with allocation-count law \(\pi\),
\[
\operatorname{BR}\ge\sum_r\pi(r)\mathcal B(r)
\ge B_{n,k_n,\varepsilon_n^{\mathrm A}}^{\dagger,\mathrm A}.
\tag{3}
\]
Taking maximum risk and then the minimax infimum proves the middle Bayes inequality in the theorem.

The exact definition gives
\[
n^{4/3}\{n^{-1}-B_{n,k_n,\varepsilon_n^{\mathrm A}}^{\dagger,\mathrm A}\}
=\max_rG_{n,k_n,\varepsilon_n^{\mathrm A}}^{\dagger,\mathrm A}(r).
\tag{4}
\]
The local, exterior, and zero-count sets partition \(\mathcal R_{n,3}\); conclusion \({\rm(A)}\) of the root and \((4)\) give the first risk inequality in \((1)\) of the statement.  Conclusion \({\rm(B)}\) gives its last two inequalities because the returned procedure is feasible in the risk game.  Thus
\[
\Lambda_{c^\dagger}-e_{\mathrm{pr},n}
\le n^{4/3}\{n^{-1}-R_n(c^\dagger)\}
\le\Lambda_{c^\dagger}+e_n^-.
\tag{5}
\]
Conclusion \({\rm(C)}\) bounds both errors by \(E(N)\) for every \(n\ge N\), proving equation \((2)\) in the statement.  Since \(E(N)\downarrow0\), rearranging \((5)\) and using \((1)\) proves the sharp expansion.  This is the algebraic corollary of the now-proved root conclusions, with no appeal to the legacy scalar prior or to an existential primal remainder.

%%% FIELD proof-cert
Fix an admissible diagonal input.  The root theorem and the normalized Airy-record lemma give a verified spline record and
\[
\alpha_n(3T_{k_n}^\dagger+4\varepsilon_n^{\mathrm A})\le1.
\tag{1}
\]
Finite-input mode with \(\delta_n=n^{-3}\) therefore terminates and gives
\[
b_n^-\le B_{n,k_n,\varepsilon_n^{\mathrm A}}^{\dagger,\mathrm A}
\le b_n^+,\qquad b_n^+-b_n^-\le n^{-3},
\tag{2}
\]
together with a total grid estimator of certified risk \(U_n\).  The exact-target van Trees lemma and orbit reduction prove
\[
b_n^-\le R_n(c^\dagger)\le U_n.
\tag{3}
\]
The all-count identity and conclusion \({\rm(A)}\) give
\[
n^{4/3}\{n^{-1}-B_{n,k_n,\varepsilon_n^{\mathrm A}}^{\dagger,\mathrm A}\}
\le\Lambda_{c^\dagger}+e_n^-.
\tag{4}
\]
Equations \((2)\)--\((4)\) imply
\[
D_n\le d_n^+\le\Lambda_{c^\dagger}+e_n^-+n^{4/3}n^{-3}
\le\Lambda_{c^\dagger}+e_n^-+n^{-1}.
\tag{5}
\]
The effective rational-grid lemma and the evaluator definition give
\[
\left[n^{4/3}(U_n-n^{-1})+U_{\Lambda,n}\right]_+
\le e_{\mathrm{pr},n},\qquad U_{\Lambda,n}\ge\Lambda_{c^\dagger}.
\]
Since \(z\le[z]_+\), this implies
\[
d_n^-\ge\Lambda_{c^\dagger}-e_{\mathrm{pr},n}.
\tag{6}
\]
Equation \((3)\) gives \(d_n^-\le D_n\le d_n^+\).  Equations \((5)\)--\((6)\) are the two asserted finite chains.

The root theorem supplies a code and proof for its computable nonincreasing rational tail \(E\).  Given rational \(\eta>0\), enumerate
\(N\ge N_{\mathrm{start}}\) and stop at the first one satisfying
\[
E(N)+N^{-1}\le\eta.
\tag{7}
\]
The certified limit of \(E\) proves termination.  For every later \(n\), equations \((5)\)--\((7)\) give all three absolute-error bounds in the statement.  On the finite inadmissible prefix, diagonal mode returns zero; since both action and target lie in \([-1,1]\), the advertised bracket
\(0\le R_n(c^\dagger)\le1\) is valid for that zero action because
\(|\tau_{c^\dagger}(y)|\le1\).  Every operation uses integers, rationals, verified rational intervals, or the ordered algebraic representation of \(n^{-1/3}\); no representation of an arbitrary real contrast is invoked.

%%% FIELD stmt-root
For \(c^\dagger=(1,-1/4,-3/4)\), there is one Turing machine extending \(\mathsf E\) with the following properties.  On input \(n\), for all sufficiently large \(n\), it uses the fixed-parameter diagonal in
\(\mathrm{def:evaluator\text{-}handle}\) and outputs an index \(k_n\), the complete verified spline record for \(\varphi_{k_n}^\dagger\), a rational
\(\varepsilon_n^{\mathrm A}>0\), rational nonnegative errors
\(e_{\mathrm{loc},n}\), \(e_{\mathrm{shell},n}\),
\(e_{\mathrm{zero},n}\), and \(e_{\mathrm{pr},n}\), and finite ordered-field certificates.  The outputs satisfy
\[
k_n\uparrow\infty,
\qquad \varepsilon_n^{\mathrm A}\downarrow0,
\qquad
\alpha_n(3T_{k_n}^\dagger+4\varepsilon_n^{\mathrm A})\le1.
\]
For
\[
\begin{aligned}
\mathcal L_n&=\left\{r\in\mathcal R_{n,3}:\min_a r_a>0,
\ \left\|n^{1/6}(r/n-q^\dagger)\right\|_\infty
\le n^{1/72}\right\},\\
\mathcal X_n&=\{r\in\mathcal R_{n,3}:\min_a r_a>0\}
\setminus\mathcal L_n,\\
\mathcal Z_n&=\{r\in\mathcal R_{n,3}:\min_a r_a=0\},
\end{aligned}
\]
the certified all-count conclusions are
\[
\begin{aligned}
G_{n,k_n,\varepsilon_n^{\mathrm A}}^{\dagger,\mathrm A}(r)
&\le\Lambda_{c^\dagger}+e_{\mathrm{loc},n}
&& (r\in\mathcal L_n),\\
G_{n,k_n,\varepsilon_n^{\mathrm A}}^{\dagger,\mathrm A}(r)
&\le\Lambda_{c^\dagger}+e_{\mathrm{shell},n}
&& (r\in\mathcal X_n),\\
G_{n,k_n,\varepsilon_n^{\mathrm A}}^{\dagger,\mathrm A}(r)
&\le\Lambda_{c^\dagger}+e_{\mathrm{zero},n}
&& (r\in\mathcal Z_n).
\end{aligned}
\tag{A}
\]
The same output contains a total finite procedure \(\widehat\tau_n^{\mathrm A}\) satisfying
\[
\max_{y\in\{0,1\}^{n\times3}}
\mathbb E\left[\{\widehat\tau_n^{\mathrm A}-
\tau_{c^\dagger}(y)\}^2\right]
\le n^{-1}-\{\Lambda_{c^\dagger}-e_{\mathrm{pr},n}\}n^{-4/3}.
\tag{B}
\]
It also outputs a code for a rational-valued nonincreasing computable function
\(E\) such that \(E(N)\to0\) and, for every \(n\ge N\),
\[
\max\{e_{\mathrm{loc},n},e_{\mathrm{shell},n},
e_{\mathrm{zero},n},e_{\mathrm{pr},n}\}\le E(N).
\tag{C}
\]
The inequalities (A)--(C), including exterior and zero-count control, are theorem conclusions and are not inputs or assumptions.  They imply
\[
\Lambda_{c^\dagger}-e_{\mathrm{pr},n}
\le n^{4/3}\{n^{-1}-R_n(c^\dagger)\}
\le\Lambda_{c^\dagger}+
\max\{e_{\mathrm{loc},n},e_{\mathrm{shell},n},
e_{\mathrm{zero},n}\},
\]
and consequently
\[
R_n(c^\dagger)=n^{-1}-\Lambda_{c^\dagger}n^{-4/3}
+o(n^{-4/3}),
\qquad
\Lambda_{c^\dagger}=C_A(5/8)^{2/3}\in(0,\infty).
\]

%%% FIELD stmt-sharp
For \(c^\dagger=(1,-1/4,-3/4)\), let \(k_n\), the verified spline record, \(\varepsilon_n^{\mathrm A}\), \(e_{\mathrm{loc},n}\), \(e_{\mathrm{shell},n}\), \(e_{\mathrm{zero},n}\), \(e_{\mathrm{pr},n}\), and \(E\) be the outputs in conclusions \({\rm(A)}\)--\({\rm(C)}\) of \(\mathrm{thm:unequal\text{-}three\text{-}arm\text{-}allocation\text{-}tube}\), and put
\[
 e_n^- = \max\{e_{\mathrm{loc},n},e_{\mathrm{shell},n},e_{\mathrm{zero},n}\}.
\]
For every sufficiently large \(n\), the corrected exact Bayes certificate and the total finite procedure \(\widehat\tau_n^{\mathrm A}\) output by that theorem satisfy
\[
 \frac1n-\{\Lambda_{c^\dagger}+e_n^-\}n^{-4/3}
 \le B_{n,k_n,\varepsilon_n^{\mathrm A}}^{\dagger,\mathrm A}
 \le R_n(c^\dagger)
 \le \max_{y\in\{0,1\}^{n\times3}}
       \mathbb E\!\left[
       \{\widehat\tau_n^{\mathrm A}-\tau_{c^\dagger}(y)\}^2
       \right]
 \le \frac1n-\{\Lambda_{c^\dagger}-e_{\mathrm{pr},n}\}n^{-4/3}.
 \tag{1}
\]
The Bayes inequality in \((1)\) is valid against every outcome-independent, possibly dependent assignment law and every randomized labeled-data estimator.  For every \(N\) beyond the finite startup of the diagonal and every \(n\ge N\),
\[
 \left|n^{4/3}\{n^{-1}-R_n(c^\dagger)\}
       -\Lambda_{c^\dagger}\right|\le E(N),
 \qquad E(N)\downarrow0.
 \tag{2}
\]
Consequently,
\[
 R_n(c^\dagger)=n^{-1}-\Lambda_{c^\dagger}n^{-4/3}
 +o(n^{-4/3}),
 \qquad
 \Lambda_{c^\dagger}=C_A(5/8)^{2/3}\in(0,\infty).
\]

%%% FIELD stmt-cert
Fix \(c=c^\dagger=(1,-1/4,-3/4)\).  Run the diagonal Turing machine furnished by the unequal-three-arm allocation-tube theorem, and write its outputs as
\(k_n\), the verified record for \(\varphi_{k_n}^\dagger\),
\(\varepsilon_n^{\mathrm A}\), \(e_{\mathrm{loc},n}\),
\(e_{\mathrm{shell},n}\), \(e_{\mathrm{zero},n}\),
\(e_{\mathrm{pr},n}\), and \(E\).  This machine terminates on every input.  Its admissibility flag is positive for every sufficiently large \(n\), and then
\[
\alpha_n\{3T_{k_n}^\dagger+4\varepsilon_n^{\mathrm A}\}\le1.
\]
Put \(\delta_n=n^{-3}\).  At every such admissible input, finite-input mode of
\(\mathsf E\) terminates and returns rationals \(b_n^-\le b_n^+\), an exact rational worst-case-risk certificate \(U_n\), and a total randomized estimator
\(\widehat\tau_n^{\mathrm A}\) under the independent allocation \(q^\dagger\), such that
\[
b_n^-\le B_{n,k_n,\varepsilon_n^{\mathrm A}}^{\dagger,\mathrm A}
\le b_n^+,
\qquad b_n^+-b_n^-\le\delta_n,
\qquad
b_n^-\le R_n(c^\dagger)\le
\max_y\mathbb E\!\left[
\{\widehat\tau_n^{\mathrm A}-\tau_{c^\dagger}(y)\}^2
\right]\le U_n.
\]
Define
\[
\begin{aligned}
e_n^-&=\max\{e_{\mathrm{loc},n},e_{\mathrm{shell},n},
e_{\mathrm{zero},n}\},\\
D_n&=n^{4/3}\{n^{-1}-R_n(c^\dagger)\},\\
d_n^-&=n^{4/3}(n^{-1}-U_n),
\qquad d_n^+=n^{4/3}(n^{-1}-b_n^-).
\end{aligned}
\]
The returned finite certificates verify
\[
\Lambda_{c^\dagger}-e_{\mathrm{pr},n}
\le d_n^-\le D_n\le\Lambda_{c^\dagger}+e_n^-,
\qquad
D_n\le d_n^+\le\Lambda_{c^\dagger}+e_n^-+n^{-1}.
\]
The same machine returns a code and proof certificate for a rational-valued nonincreasing computable function \(E\) satisfying \(E(N)\to0\) and
\[
\max\{e_{\mathrm{loc},n},e_{\mathrm{shell},n},
e_{\mathrm{zero},n},e_{\mathrm{pr},n}\}\le E(N)
\qquad(n\ge N).
\]
Consequently, for every rational \(\eta>0\), a terminating search returns a computable
\(N^\dagger(\eta)\) such that, for all \(n\ge N^\dagger(\eta)\),
\[
d_n^-\le D_n\le d_n^+,
\qquad
|d_n^--\Lambda_{c^\dagger}|\le\eta,
\qquad
|D_n-\Lambda_{c^\dagger}|\le\eta,
\qquad
|d_n^+-\Lambda_{c^\dagger}|\le\eta.
\]
On the finite inadmissible prefix, the evaluator returns its flag and the deterministic zero estimator with the valid bracket
\(0\le R_n(c^\dagger)\le1\); hence the finite-data procedure is total for every \(n\).  No computability representation of an arbitrary real contrast is used.
